\documentclass[11pt,a4paper]{article}
\usepackage[top=2cm,bottom=2cm,left=2.2cm,right=2.2cm]{geometry}
\usepackage[utf8]{inputenc}
\usepackage[T1]{fontenc}
\usepackage{mathptmx}
\usepackage{amssymb}
\usepackage{xcolor}
\definecolor{acc}{HTML}{00723F}
\definecolor{gr}{HTML}{555555}
\usepackage{enumitem}
\setlist[itemize]{leftmargin=1.3em,itemsep=1pt,topsep=2pt,parsep=0pt}
\usepackage{titlesec}
\titleformat{\section}{\large\bfseries\color{acc}}{\thesection}{0.6em}{}
\titleformat{\subsection}{\normalsize\bfseries}{\thesubsection}{0.6em}{}
\titlespacing*{\section}{0pt}{14pt}{5pt}
\titlespacing*{\subsection}{0pt}{9pt}{3pt}
\usepackage{array}
\usepackage{longtable}
\usepackage{colortbl}
\newcommand{\job}[1]{\par\vspace{2pt}\noindent\textcolor{gr}{\textit{Job: #1}}\par\vspace{2pt}}
\newcommand{\turn}[1]{\par\vspace{3pt}\noindent\textbf{The point to land:} #1\par}
\newcommand{\len}[1]{\hfill\textcolor{gr}{\small[#1]}}
\newcommand{\todo}{$\square$~}
\setlength{\parindent}{0pt}
\pagestyle{plain}

\begin{document}

\begin{center}
{\LARGE\bfseries\color{acc} Chapter 1 --- Writing Notes}\\[3pt]
{\small Project Framework, State of the Art, and Needs Analysis}\\[2pt]
{\small\color{gr} Points to make and facts to use. Write the sentences yourself.}
\end{center}

\vspace{6pt}
\hrule
\vspace{8pt}

\textbf{How to use this:} read a section's notes, close the page, write. Drafting
with notes in front of you produces reshaped notes, not your own writing. Write
badly first and fix afterwards --- the blank page costs more than the editing.
One idea per paragraph. State the point, then the evidence, not the reverse.

\vspace{4pt}
Target for the whole chapter: 12--15 pages.

\vspace{8pt}
\hrule

% =====================================================================
\section*{Introduction to the chapter \len{$\sim$½ page, 2 paragraphs}}
\job{tell the reader what this chapter does and why it matters}

\textbf{Para 1}
\begin{itemize}
\item where the project came from, what it was asked to fix
\item names: the host institution, the client, the document
\item flag that most of the chapter is about how the document gets used today
\end{itemize}

\textbf{Para 2}
\begin{itemize}
\item the requirements at the end weren't decided up front
\item they're what's left after watching the current process fail
\item $\rightarrow$ that's why failures come first, requirements after
\item closes with: existing tools, why none fit, what was proposed
\end{itemize}

\turn{this chapter isn't a survey --- it's where the specification comes from}

% =====================================================================
\section*{1.1 General Framework of the Project \len{$\sim$½ page + table}}
\job{the administrative framing, quickly}

\begin{itemize}
\item final-year project, Master's in Data Sciences, Polytech Intl
\item six-month internship; $\sim$4 months on the build
\item rest: framing the subject, learning the source material, writing this
\end{itemize}

\textbf{Then 1--2 sentences: what kept you on it.} Three angles, pick one:
\begin{itemize}
\item[(a)] two problems pulling opposite ways --- extraction rewards rules you can
check against a page; language models give flexibility with no guarantee of faithfulness
\item[(b)] a page a human reads easily breaks every library you point at it ---
finding out why meant learning to read the file the way the printer wrote it
\item[(c)] the interesting problem wasn't building a chatbot --- it was working out
what a chatbot must \emph{not} be allowed to do when the subject is who controls money
\end{itemize}

\textbf{\todo Table 1.1 --- Framework of the project.} Rows: project type, host
institution, programme, academic year, duration, client organisation, academic
supervisor, professional supervisors, source document, scope processed.

% =====================================================================
\section*{1.2 Host Organisation \len{$\sim$1½ pages}}

\subsection*{1.2.1 Presentation of Polytech Intl}
\job{establish the institution; keep it factual, not promotional}

\begin{itemize}
\item \'Ecole Polytechnique Internationale Priv\'ee de Tunis
\item founded 2013 \quad\textcolor{gr}{\small(verify on the school's own site)}
\item Rue du Lac d'Annecy, Les Berges du Lac 1, Tunis
\item founded largely by institutional figures from \textbf{banking, insurance, industry}
\item programmes: computer science (big data, business intelligence, software
engineering, IT-finance, embedded systems), mechatronics, industrial engineering,
financial engineering, telecommunications, architecture
\item international mobility --- partner schools in Europe, Asia, North America
\item member of the Agence Universitaire de la Francophonie \textcolor{gr}{\small(verify)}
\item teaching built around projects + compulsory final internship
\end{itemize}

\textbf{Then connect to your project:}
\begin{itemize}
\item Data Sciences Master's spans stats/ML, data engineering, and the software
side of making an analysis actually run
\item final project expected to show all three
\item $\rightarrow$ that's why this went to a deployed application, not a clean dataset
\end{itemize}

\subsection*{1.2.2 Project Context and the Client Organisation}
\job{who the client is, and why an institution like that needs a DAM at all}

\textbf{Positioning --- get this exactly right:}
\begin{itemize}
\item Polytech Intl \textbf{hosted and supervised}
\item AfDB was the \textbf{client}: supplied the source document and the requirements
\item do \emph{not} describe yourself as working at or within AfDB anywhere
\end{itemize}

\textbf{AfDB facts:}
\begin{itemize}
\item founded 1964, HQ Abidjan, C\^ote d'Ivoire
\item 81 member countries --- 54 African, 27 non-regional
\item purpose: economic development and social progress across regional member countries
\item Group = the Bank + African Development Fund (concessional lending) + Nigeria Trust Fund
\item RDGN: Tunisia, Algeria, Morocco, Mauritania, Libya --- hub in Tunis
\textcolor{gr}{\small(confirm official wording with Yasser or Surya)}
\end{itemize}

\textbf{\todo One sentence only you can write:} how the subject reached you ---
who proposed it, how they described the problem.

\textbf{Optional nice line:} a school founded partly by people from banking,
supervising a project built for a bank.

\textbf{Closing paragraph --- why a real client changed the work:}
\begin{itemize}
\item material couldn't be tidied to suit the parser
\item correctness = what the printed pages say, not a convenient metric
\item every extraction claim in the report is checkable against its page
\item $\rightarrow$ that constraint runs through Chapter 2
\end{itemize}

\textbf{\todo Table 1.2 --- The client at a glance.} Rows: institution, founded,
HQ, membership, purpose, Group entities, regional office concerned, regional hub.

\newpage
% =====================================================================
\section*{1.3 Study of the Existing System and Critical Analysis \len{$\sim$5 pages}}

\subsection*{1.3.1.1 Definition and stakes of the DAM \len{2--3 paragraphs}}
\begin{itemize}
\item delegation framework answers: for a given institutional act, who may do which part
\item expressed as a matrix
\item rows = activities, hierarchical: chapters $\rightarrow$ processes $\rightarrow$ tasks
$\rightarrow$ child tasks $\rightarrow$ threshold variants (keyed by money band)
\item columns = roles: country office staff $\rightarrow$ sector managers $\rightarrow$
directors $\rightarrow$ VPs $\rightarrow$ Board
\item cell = code saying what that role does for that activity
\end{itemize}

\textbf{\todo Table --- authority notation.} I / C1,C2 / C3,C4 / R / A / (i)

\textbf{Two traps in the notation --- both matter later:}
\begin{itemize}
\item ``C'' covers two unrelated ideas split only by the digit:
C1,C2 = check and verify; C3,C4 = consult (Board / independent function)
\item the informed marker isn't a letter at all --- it's parenthesised, and
behaves differently at every processing stage
\end{itemize}

\subsection*{1.3.1.2 Governance and compliance requirements \len{1--2 paragraphs}}
\begin{itemize}
\item the DAM isn't advisory --- it's the reference internal audit checks against
\item exists partly because public and multilateral funding brings governance expectations
\end{itemize}
Three consequences for any system built on it:
\begin{itemize}
\item[1.] a wrong answer surfaces at audit, not during the work
\item[2.] completeness = correctness: naming the approver but omitting a mandatory
verifier is incomplete in a way that misleads, though nothing in it is false
\item[3.] traceability: an answer you can't verify gives you no basis for relying on it
\end{itemize}
\turn{consequence 2 is where your FR4 comes from --- flag it}

\subsection*{1.3.1.3 Challenges in consulting it in practice \len{2--3 paragraphs}}
\begin{itemize}
\item tables are wide $\rightarrow$ role headers rotated 90° to fit the page
\item reader must map a sideways label to a code centimetres below it
\item and get the adjacent column right too --- neighbouring roles often differ on the same activity
\item superscript footnote digits sit against the codes; telling them apart requires knowing which is which
\item some rows hold no responsibilities at all --- just a cross-reference, so you
restart in a different chapter
\item new staff need the document most and are least equipped to navigate it
\end{itemize}
\turn{consequence: people ask a colleague instead, or copy what was done last time.
Neither is auditable. Both propagate errors quietly.}

\subsection*{1.3.2 Study of Existing Systems}
\textit{Opening: four categories exist, none works, each fails for its own reason,
and those reasons become the spec.}

\textbf{1.3.2.1 Manual consultation and document search \len{2 paragraphs}}
\begin{itemize}
\item today: open PDF, Ctrl+F
\item full-text search = character string matching
\item works if you know the exact title or the identifier; close to useless otherwise
\item ``approve'' $\rightarrow$ hundreds of hits, page order, no relevance ranking
\item half-remembered name $\rightarrow$ nothing unless wording matches exactly
\item \textbf{deeper point:} search finds \emph{places}; you wanted an \emph{answer}
\item it says page 24; it can't say who approves, because that needs: page holds a
table + headers rotated + letter in cell = a role's authority
\end{itemize}

\textbf{1.3.2.2 Generic document QA assistants \len{2--3 paragraphs}}
\begin{itemize}
\item named: ChatPDF, NotebookLM, Microsoft Copilot, file-upload chat modes
\item mechanism (state it accurately --- you'll be asked): split into passages
$\rightarrow$ embed as vectors $\rightarrow$ retrieve nearest to question
$\rightarrow$ model writes from those. This is RAG; use the term.
\item works on prose; fails here
\item failure is \textbf{structural, not model quality} --- don't let a reader think
a better model fixes it
\item flatten page $\rightarrow$ spatial relationships destroyed
\item rotated header becomes a loose string attached to nothing
\item retrieved passage can hold the activity name AND the letter A, having lost
whose column the A was in
\item answer is gone before the model sees anything
\end{itemize}
\textbf{\todo Worth 15 minutes:} upload 2--3 DAM pages to ChatPDF or NotebookLM,
ask ``who approves 2.126'', screenshot it. Wrong or vague $\rightarrow$ becomes a
figure, and turns this subsection from argument into evidence.

\textbf{1.3.2.3 Commercial document-extraction services \len{2 paragraphs}}
\begin{itemize}
\item named: Azure Document Intelligence, AWS Textract, Google Document AI, LlamaParse
\item be fair: markedly better at tables than a general library; the serious alternative
to hand-written rules
\end{itemize}
Three reasons ruled out:
\begin{itemize}
\item[1.] per-page cost + sending an internal governance document to a third party
--- an institutional decision, not a student's
\item[2.] tuned for invoices, receipts, contracts, forms --- not a bespoke authority
matrix with rotated headers and superscript-qualified codes
\item[3.] \textbf{no account of why a cell came out wrong.} A rule from a measured
character offset is checkable against its page. A confidence score isn't.
\end{itemize}
\turn{reason 3 is your strongest and connects to the thesis's whole argument --- give it the most room}

\textbf{1.3.2.4 Critical analysis \len{2--3 paragraphs + table}}
\begin{itemize}
\item none fail from poor engineering --- each fails on an assumption that doesn't hold
\item search assumes $\rightarrow$ user can phrase it in the document's vocabulary
\item assistants assume $\rightarrow$ flattening preserves meaning (false for positional tables)
\item extraction assumes $\rightarrow$ document resembles their training forms
\item GRC platforms (4th category) assume $\rightarrow$ institution re-enters its rules
into their data model = change programme, not a tool
\end{itemize}
\textbf{The shared failure --- put weight here:}
\begin{itemize}
\item none has any concept of being wrong
\item search returns results, model returns text; neither separates
``answer I can support'' from ``answer I can't''
\item for a compliance reference that distinction is the whole point
\end{itemize}

\textbf{\todo Comparison table.} Columns: Search $\cdot$ Assistants $\cdot$
Extraction $\cdot$ GRC $\cdot$ This work. Rows: preserves positional meaning;
answers vs locates; traceable to a source row; refuses when it can't answer;
surfaces mandatory obligations; tolerates typos; document stays in-house; no
per-page cost; deployable here.

\textcolor{gr}{\small Don't mark a competitor ``No'' on something you haven't tested
--- use Partly or N/A and you're safe if asked.}

\turn{read down your own column --- that list \emph{is} the specification. Say so, then point at 1.5.}

\newpage
% =====================================================================
\section*{1.4 Functional Needs Analysis \len{$\sim$3 pages, mostly tables}}

\subsection*{1.4.1 Functional requirements}
Short intro para: these weren't invented --- each corresponds to something a user
currently does by hand, or a failure observed in 1.3.

\textbf{\todo Table.} Core (FR1--FR10):
\begin{itemize}
\item FR1 direct identifier lookup --- full confidence, no ambiguity
\item FR2 natural-language activity lookup + reports its confidence
\item FR3 authority-role queries (approve / review / check / initiate / inform)
\item FR4 mandatory obligation surfacing --- check/verify + informed parties, whatever was asked
\item FR5 cross-reference resolution --- follow the pointer automatically
\item FR6 typo and phrasing tolerance
\item FR7 honest refusal --- unknown id, out of scope
\item FR8 conversational continuity --- follow-ups without restating the subject
\item FR9 structural analytics --- counts, code distribution, top roles, per-chapter
\item FR10 multilingual questions --- FR/ES/PT/AR, without translating role names or ids
\end{itemize}
Second phase (FR11--FR16): authenticated access; voice input; document attachment
with a distinct trust boundary; fully localised interface incl. RTL; conversation
export as PDF with provenance; read-only sharing by QR code.

\subsection*{1.4.2 Non-functional requirements}
\textbf{\todo Table.}
\begin{itemize}
\item NFR1 factual fidelity --- no role/code/qualification not in the source.
\textbf{Takes precedence over everything else.} Say that explicitly.
\item NFR2 traceability --- every answer attributable to a specific record
\item NFR3 availability without external services --- degrade to templates, don't fail
\item NFR4 response time --- seconds on free hosting $\rightarrow$ rules out re-parsing per request
\item NFR5 reproducibility --- same facts every run, independent of model sampling
\item NFR6 portability --- identical on laptop and host
\item NFR7 usable on a phone
\item NFR8 credential protection --- no recoverable passwords, tokens expire
\item NFR9 upload safety --- type and size limits, never executed
\end{itemize}

\subsection*{1.4.3 Project constraints}
\textbf{\todo Table.} Constraint $\rightarrow$ consequence:
\begin{itemize}
\item no labelled training data $\rightarrow$ ruled out supervised model; pushed to lexical retrieval
\item document scope: 250+ pages, 79 processed $\rightarrow$ rest as future work
\item no infrastructure budget $\rightarrow$ free tier $\rightarrow$ forced build-time caching
\item no ground truth $\rightarrow$ correctness established against real pages $\rightarrow$
set the pace of the whole data understanding phase
\end{itemize}

% =====================================================================
\section*{1.5 Proposed Solution \len{$\sim$2 pages}}

\subsection*{1.5.1 General presentation}
\begin{itemize}
\item the core idea: separate \emph{what the facts are} from \emph{how they're said}
\item facts $\leftarrow$ structured knowledge base, parsed once offline, rules derived
from the document's own geometry, validated against printed pages
\item phrasing $\leftarrow$ language model, given facts already retrieved, restricted to rewording
\item between them: a verification step
\item any answer that dropped or altered a role name is rejected $\rightarrow$
deterministic template shown instead
\item user still gets a correct answer, just less conversational
\end{itemize}
\turn{this is the mechanism that makes a generative component usable at all here}

\subsection*{1.5.2 Global system architecture}
Four layers:
\begin{itemize}
\item[1.] offline extraction --- PDF $\rightarrow$ validated node set, cached, shipped with the app
\item[2.] knowledge --- node set as a directed graph + lookup and search over it
\item[3.] agent --- interprets questions, resolves to activities, builds deterministic
answers, applies the model + verification
\item[4.] presentation --- REST interface + authenticated web client (chat, dashboard,
attachment, export, sharing)
\end{itemize}
\textbf{\todo Figure: fig1\_architecture.png}

\textbf{Worth a sentence:} extraction happens offline and is cached, so the running
app never opens the PDF. It loads a prepared node set at startup. That was a
deployment decision --- Chapter 5 covers what happened when it wasn't.

\subsection*{1.5.3 Selected technologies}
\textbf{\todo Table.} Technology $\rightarrow$ role. Python; pdfplumber (character-level
coordinates --- required, not convenient); Pydantic (schema validation at construction);
networkx (directed multigraph --- a role can hold several responsibilities on one
activity); scikit-learn (TF-IDF + cosine); FastAPI; React + Vite; Groq + Ollama
(behind one interface, so either or neither); Web Speech API; python-jose + passlib;
ReportLab; statsmodels; Docker; Render.

\textcolor{gr}{\small For each, give the reason not just the name. ``pdfplumber, because
the superscript rule needs per-character vertical offset and font size'' beats
``pdfplumber for PDF extraction''.}

% =====================================================================
\section*{Conclusion to the chapter \len{$\sim$½ page, 1--2 paragraphs}}
\begin{itemize}
\item the DAM answers an important question; the way it's published makes the answer expensive
\item manual consultation is slow and error-prone because the meaning lives in the layout
\item generic assistants fail for the same reason --- with the added risk that they fail \emph{fluently}
\item the requirements above respond to both
\item the architecture separates retrieval from phrasing so a model's flexibility
can be used without inheriting its unreliability
\item hand off: next chapter examines the source document and what makes extracting it hard
\end{itemize}

\vspace{10pt}
\hrule
\vspace{6pt}
\textbf{Your numbers, for cross-reference:} 250+ page document, 79-page subset ---
327 records (3 chapters, 35 processes, 138 tasks, 102 child tasks, 49 threshold
variants) --- 1,776 responsibility entries --- 131 roles --- 0.56\% unresolved ---
5.48 responsibilities per record --- 29.63\% with none --- 14 cross-references ---
graph 459 nodes / 2,108 edges --- 374 tests across 23 files.

\end{document}
